\section{Results}
\label{sec:results}

Our experiments confirm the four-step causal mechanism and quantify the calibration-alignment divergence curve with high statistical confidence in the primary dataset and medium confidence in the independent replication.

\subsection{Step 1: Signal Co-existence (H-E1)}

Both RM score and gold human preference co-exist as separable, non-constant series. In Coste et al.\ data, RM score variance $= 1.96$ and gold variance $= 0.25$ across 10 KL checkpoints; in Gao et al.\ data, rm\_var $= 2.42$ and gold\_var $= 0.31$ across 11 KL levels available for signal characterization (note: one boundary checkpoint at KL $= 0$ is included here for variance estimation but excluded from the regression in Step 4---see Section~\ref{subsec:replication}---yielding $n=10$ paired observations for H-M4). Both signals are non-constant and separable.

\textbf{What this means:} The measurement infrastructure for bidirectional alignment evaluation already exists in published RLHF experiments---both proxy and gold signals are recoverable and distinct. Figure~\ref{fig:dual_coste} (\texttt{dual\_axis\_coste2023.png}) illustrates both trajectories.

\subsection{Step 2: Directional Divergence (H-M1)}

The reward model score rises monotonically (Spearman $\rho(\text{KL}, \text{RM}) = 1.000$, $p < 0.0001$) while gold preference peaks at KL $\approx 2.0$ nats (gold $= 0.63$) and declines to $0.38$ by KL $= 8.0$ nats---a 40\% decline from peak. Final proxy-gold divergence $= 1.70$.

Figure~\ref{fig:trajectory} (\texttt{trajectory\_dual\_axis.png}) shows the diverging trajectories. \textbf{What this means:} Sustained RLHF optimization does not plateau human preference---it reverses it. The proxy and human judgment signals move in opposite directions once KL $> 2$ nats.

\textbf{Caveat:} $\rho = 1.000$ is unexpectedly perfect for real experimental data; we attribute this most likely to digitization idealization. Gao et al.\ $R^2 = 0.701$ (Section~\ref{subsec:replication}) is more representative of real experimental noise.

\subsection{Step 3: Gap Positivity and Monotonicity (H-M2)}

The normalized divergence gap is strictly positive at all five high-KL checkpoints (KL $> 3.5$ nats), with $\rho(\text{gap}, \text{KL}) = 1.000$ within the high-KL subset.

\begin{table}[h]
\centering
\caption{Normalized divergence gap at high-KL checkpoints.}
\label{tab:gap}
\begin{tabular}{ll}
\toprule
KL budget (nats) & gap $= \text{RM\_norm} - \text{gold\_preference}$ \\
\midrule
4.0 & $+0.277$ \\
5.0 & $+0.388$ \\
6.0 & $+0.484$ \\
7.0 & $+0.560$ \\
8.0 & $\mathbf{+0.620}$ \\
\bottomrule
\end{tabular}
\end{table}

Maximum gap $= 0.620$---the proxy score exceeds gold preference by more than half the normalized scale. Figure~\ref{fig:gap} (\texttt{gap\_curve.png}) shows the gap trajectory. \textbf{What this means:} Once KL $> 3.5$ nats, the calibration-alignment divergence is systematic and monotonically growing---not a transient artifact.

\subsection{Step 4: Divergence Curve Quantification (H-M3 --- Primary)}

\begin{table}[h]
\centering
\caption{OLS regression results, Coste et al.~\cite{Coste2023Reward} (primary dataset).}
\label{tab:regression_coste}
\begin{tabular}{ll}
\toprule
Metric & Value \\
\midrule
Slope $\beta$ & $\mathbf{0.1433\ \text{nat}^{-1}}$ \\
Intercept $\alpha$ & $-0.4016$ \\
$R^2$ & $\mathbf{0.9577}$ \\
$p$-value (Wald $t$-test) & $\mathbf{8.89 \times 10^{-7}}$ \\
$t$-statistic & $13.461$ \\
$N$ & $10$ \\
Parametric 95\% CI & $[0.119, 0.168]$ \\
Bootstrap 95\% CI ($N=10{,}000$) & $[0.117, 0.177]$ \\
\bottomrule
\end{tabular}
\end{table}

All three gate conditions satisfied: $\beta > 0$, $p < 0.05$, $R^2 > 0.5$. Both CIs are strictly positive---HIGH confidence for P1.

Figure~\ref{fig:regression} (\texttt{regression\_scatter.png}) shows the scatter plot with OLS line and 95\% CI. \textbf{What this means:} A single linear model explains 96\% of the variance in the calibration-alignment divergence gap across 10 RLHF checkpoints. Each additional nat of KL budget adds 0.143 normalized units to the divergence.

\subsection{Step 4 (Replication): Cross-Dataset Replication (H-M4 --- Replication)}
\label{subsec:replication}

\begin{table}[h]
\centering
\caption{Cross-dataset comparison of regression results.}
\label{tab:crossdataset}
\begin{tabular}{lll}
\toprule
Metric & Coste et al.~\cite{Coste2023Reward} & Gao et al.~\cite{Gao2023Scaling} \\
\midrule
$\beta$ (slope) & $0.1433\ \text{nat}^{-1}$ & $\mathbf{0.1599\ \text{nat}^{-1}}$ \\
$R^2$ & $0.9577$ & $0.7008$ \\
$p$-value & $8.89 \times 10^{-7}$ & $\mathbf{0.0025}$ \\
Parametric 95\% CI & $[0.119, 0.168]$ & $[0.075, 0.245]$ \\
Bootstrap 95\% CI & $[0.117, 0.177]$ & $[-0.020, 0.236]$ \\
Confidence & HIGH & MEDIUM \\
\bottomrule
\end{tabular}
\end{table}

Cross-dataset slope ratio $\beta_\text{Gao} / \beta_\text{Coste} = \mathbf{1.116}$---slopes within 12\% across independent model families.

Figure~\ref{fig:slopes} (\texttt{fig3\_cross\_dataset\_slopes.png}) shows the slopes with error bars. Figure~\ref{fig:overlay} (\texttt{fig4\_dual\_overlay.png}) shows both datasets' regression lines.

\textbf{Note on Gao checkpoint count:} The dual-axis trajectory plot (Appendix~\ref{app:gao}, Figure~\ref{fig:gao_dual}) displays all 11 available KL levels from Gao et al.\ for signal visualization. The regression (H-M4) uses $n=10$ paired observations, excluding the KL$=0$ boundary point where RM score is at its minimum anchor value and the gap is not yet informative as a regression predictor. This exclusion is consistent with the normalization procedure and is reflected in the datasets table (Table~\ref{tab:datasets}).

\textbf{Caveat (Gao replication):} Bootstrap CI $[-0.020, 0.236]$ marginally overlaps zero due to the non-monotone gap shape at low KL in Gao data (gap is initially negative before crossover at $\sim$3.5 nats). The parametric Wald $t$-test ($p = 0.003$) and parametric CI $[0.075, 0.245]$---entirely positive---serve as the primary replication criterion, as pre-registered.

\subsection{Summary}

\begin{table}[h]
\centering
\caption{Hypothesis validation summary.}
\label{tab:summary}
\begin{tabular}{lllp{5.5cm}}
\toprule
Hypothesis & Gate & Status & Key Result \\
\midrule
H-E1 (Co-existence) & MUST\_WORK & PASS & 10 paired KL levels; both signals separable \\
H-M1 (Mechanism) & MUST\_WORK & PASS & $\rho(\text{KL},\text{RM})=1.000$; gold $-40\%$ from peak \\
H-M2 (Gap positivity) & MUST\_WORK & PASS & 5/5 high-KL positive; max\_gap$=0.620$ \\
H-M3 (Coste slope) & MUST\_WORK & PASS & $\beta=0.143$, $p=8.89\times10^{-7}$, $R^2=0.958$ \\
H-M4 (Gao replication) & SHOULD\_WORK & PASS & $\beta=0.160$, $p=0.003$; ratio$=1.116$ \\
\bottomrule
\end{tabular}
\end{table}
